% =====================================================================
\section{Deployment Experiences and Lessons}
\label{sec:deployment}
% =====================================================================

\subsection{Use Cases of Grounding, ML, and DL}

\begin{table*}[t]
\caption{Deployment case studies: where each technology family is exercised and
what it produced.}
\label{tab:cases}
\small
\begin{tabular}{@{}p{2.8cm}p{2.0cm}p{5.2cm}p{5.6cm}@{}}
\toprule
Case & Family & What was done & Core impact \\
\midrule
Corpus-scale gate evaluation & Grounding & 5{,}200 synthetic résumés, roughly
100{,}000 adversarial trials, run offline against real pipeline modules &
Surfaced a silent 0.82\% false-rejection defect invisible to the unit tests \\
Threshold calibration & DL & Equal-error sweep on LFW, 2{,}200 train and 1{,}000
held-out pairs & False rejection reduced from 0.414 to 0.034 \\
Nine-model comparison & ML & Ten grouped folds, Friedman plus Holm-corrected
Wilcoxon, synthetic positive control & Established that algorithm choice is not
the bottleneck on this task \\
Within-query diagnostic & ML & 4{,}112 within-résumé pairs, entity effects
cancelled by construction & Separated matching from candidate ranking; refuted
our own hypothesis \\
Sufficiency pipeline validation & ML & Planted-weight synthetic generator with
two weightless noise features & Pipeline recovered the planted structure and
correctly ignored the noise features \\
Independent security audit & All & Structured review of 17 public and 23
authenticated endpoints, 3 upload paths & 10 findings (5 high, 5 medium); 9
fixed, 1 mitigated \\
\bottomrule
\end{tabular}
\end{table*}

\subsection{Lessons Learned from Real Deployments}

The most instructive material came from auditing a reference implementation of a
comparable system, read in full before deciding what this project should inherit.
Its computer-vision core was sound and worth porting. Its decision layer, however, was largely fabricated, and the fabrications shared a single root cause.

\paragraph{A detector that never ran, and a heuristic that filled the gap.}
The object detector was configured with an invalid model name, so it threw on
every load and real detection never executed once. What filled the gap was a
luminance heuristic: if the left or right third of the frame deviated more than
60\% from mean brightness for three consecutive frames, it logged a phone or book detection, labelled \emph{camera verified}. A window, a desk lamp, or a light-coloured shirt was sufficient, and those fabricated events then fed an escalating penalty system. Compounding it, the load-failure handler set a status
string reading \emph{proctoring engine ready}, and that status was never rendered
anywhere, so a total detector outage was invisible.

\paragraph{A confident pass manufactured from missing data.}
The identity endpoint returned a verified result with a similarity of 98.7 and a
confidence of 0.98 when no enrolled face profile existed at all. The absence of a record produced a high-confidence pass, not an error.

\paragraph{Analysis that measured nothing.}
A résumé analyser substring-matched a fixed keyword list and, on zero matches,
returned a plausible-looking hardcoded skill list; its score was floored so an
unrelated résumé could never fall below a fifth of the scale. A speech-analytics
component took text rather than audio and derived fluency metrics from it.

\begin{table}[t]
\caption{Lessons carried into this system's architecture.}
\label{tab:lessons}
\small
\begin{tabular}{@{}p{2.1cm}p{5.3cm}@{}}
\toprule
Lesson & Architectural response \\
\midrule
A fallback that fabricates is worse than an outage & A failed measurement has its
own type carrying no value; nothing can substitute for it \\
Silent failure is the dangerous failure & Component health is surfaced where a
human sees it; corpus-scale evaluation with an explicit true-positive metric \\
Plausible output invites trust & Outputs are traceable to source spans, so
plausibility is checkable rather than persuasive \\
Hidden weights become verdicts & The composite score is removed, not improved \\
Reasoned is not correct & Security-adjacent constants require calibration studies
with reported error rates \\
\bottomrule
\end{tabular}
\end{table}

\begin{table}[t]
\caption{Where deep learning is and is not used, and why.}
\label{tab:dl-domains}
\small
\begin{tabular}{@{}p{2.0cm}p{5.4cm}@{}}
\toprule
Domain & Decision \\
\midrule
Face detection and embedding & Used, pretrained, frozen; genuinely a perception problem \\
Semantic comparison & Used for clustering and relational features; similarity is
not treated as entailment \\
Claim proposal & Used, but gated; the model never authors \\
Affect or emotion inference & \textbf{Not used.} No defensible accuracy on a
diverse candidate pool, and it maps onto disability, neurodivergence, and
cultural difference in expression \\
Demographic attributes & \textbf{Never loaded.} The bundled classifier is
excluded by an allow-list \\
Deception or cheating estimation & \textbf{Not built.} A learned cheating score
has no ground truth and cannot be calibrated \\
\bottomrule
\end{tabular}
\end{table}

Table~\ref{tab:dl-domains} records an argument worth stating explicitly. A
proposed live dashboard included confidence, stress, and eye-contact percentages
feeding back into question selection. We did not wire them in, and the reasoning
generalises: a stress percentage that silently changes which questions a
candidate is asked \emph{is} a hidden score, and one the candidate can neither
see nor contest. Every adaptive behaviour that dashboard promised is reachable from what the candidate \emph{said} --- observable, quotable, and defensible in a rejection conversation.

\subsection{Honest Position on Maturity}

Three gaps are stated openly here. Identity matching is implemented and
unit-tested but is currently only a presence gate on the production candidate
path, so the system does \emph{not} today perform continuous identity
verification. The face threshold is calibrated on a public benchmark rather than
on candidates. No structured recruiter validation interviews have been conducted,
which makes the premise that recruiters prefer a claim audit to a score the
project's largest untested assumption. The résumé--posting fit model reached
0.657 AUROC in earlier work and is \emph{not deployed}; no runtime module imports
it, and \S\ref{sec:ml} explains why that number is weak.

% =====================================================================
\section{Future Directions and Research Opportunities}
\label{sec:future}
% =====================================================================

\begin{table*}[t]
\caption{Research areas in evidence-based screening: core challenges and key
questions.}
\label{tab:research-areas}
\small
\begin{tabular}{@{}p{3.0cm}p{5.6cm}p{6.4cm}@{}}
\toprule
Research area & Core challenge & Key question \\
\midrule
Deployment-aware audit architectures & Continuous provenance without continuous
surveillance & How much identity assurance is obtainable from intermittent,
consented checks with explicit unmeasured intervals? \\
Explainable and contestable screening & Explanations that support appeal, not
just interpretation & Does a claim audit change human decisions relative to a
score, measured by inter-reviewer agreement? \\
Privacy-preserving process evidence & Keystroke and timing data are highly
identifying & Can within-session process deviation be measured on-device so only
derived scores leave the client? \\
Within-query evaluation standards & Pooled metrics conflate capabilities & Should
conditioning on the recurring entity become a reporting requirement for paired
benchmarks? \\
Benchmark provenance & Undocumented corpora beget undocumented derivatives &
What minimum provenance metadata should a hiring benchmark carry before it is
usable for published comparison? \\
Calibrated abstention & Choosing an operating point for ``unknown'' & What
abstention rate maximises usefulness while preserving honesty? \\
\bottomrule
\end{tabular}
\end{table*}

\subsection{Deployment-Aware Audit Architectures}

The immediate engineering task is wiring the calibrated matcher into the live
candidate session, which \S\ref{sec:deployment} records as not yet done. The
research question underneath it is more interesting: continuous verification and
candidate privacy pull in opposite directions, and the honest resolution is
probably not more surveillance but better accounting for intervals in which
nothing was measured.

\subsection{Explainable and Contestable Screening}

The project's core premise is untested: that a claim audit is more useful to a reviewer than a score. A controlled comparison of reviewers given an audit against reviewers given a score, measured by inter-reviewer agreement on identical sessions, would test it directly. We regard this as the single most valuable study the project
does not yet have, and note that it could refute the premise.

\subsection{Privacy-Preserving Process Evidence}

Process signals are informative and dangerous. Our position is that a deviation
in how an answer was produced should select a \emph{question}, never raise a \emph{flag}. The candidate's answer is the evidence. Whether that position survives contact with real candidates is, as yet, untested. Technically, the promising direction is
computing deviation on-device so that raw timings never leave the client and only a small number of derived scores are transmitted. That is a privacy property, not an optimisation.

\subsection{Within-Query Evaluation Standards}

\S\ref{sec:withinquery} showed that a representation change worth $+0.041$ pooled
AUROC delivered no measurable within-query improvement. The diagnostic costs one
pass over predictions already computed, needs no additional labels, and has a
fixed null of 0.5. We suggest reporting it alongside pooled metrics as a matter
of course for any paired benchmark with a recurring entity, and treating a
divergence between the two as information rather than as an anomaly to explain
away.

% =====================================================================
\section{Conclusion}
\label{sec:conclusion}
% =====================================================================

CogniHire replaces the composite hiring score with a per-claim audit. Claims are
extracted under a grounding constraint enforced by an architectural boundary
rather than by prompt discipline, probed in a live interview, and reported as verdicts traceable to evidence, including explicit statements about what was not measured. A human decides.

The methodological contribution is the repeated conversion of policies into
structures. ``The AI must not author claims'' becomes an import boundary with a
failing test. ``Unmeasured is not passed'' becomes a type with no field to hold a
fabricated value. ``No demographic inference'' becomes a module that is never
loaded. Each is verifiable by inspection in seconds, and none of them decays as attention moves elsewhere.

Three empirical findings stand out. The grounding gate rejected 100\% of
paraphrases and negation traps across roughly 100{,}000 adversarial trials while
admitting 100\% of verbatim claims, making ``the model selects, never authors'' a measured property, not a design aspiration. A carefully reasoned
similarity constant of 0.50 would have falsely rejected 41.4\% of genuine
identity matches; measurement reduced that to 3.4\%. And nine classifiers
spanning six families proved statistically indistinguishable at the top on a
public screening benchmark, while a within-query diagnostic showed that the
pooled metric's apparent gains came from ranking candidates rather than matching
them to roles.

That last finding deserves the final word, because it began as a different one.
The evidence initially indicated that the benchmark's labels were not relational
at all, and a direct test refuted it: four-fifths of the label variance is
within-résumé and models recover a real part of it. What survived was narrower: a decoupling between leaderboard progress and task progress that no pooled metric can reveal. We report the sequence rather than the conclusion alone,
because a diagnostic capable of distinguishing two hypotheses is only shown to
have that power when it actually overturns one.

We report a weak result and a rejected component alongside strong ones. A project
that publishes only its successes offers no evidence that it evaluates itself
honestly, and honest evaluation is what this system claims to provide to the
people it screens.

\section*{Data and Code Availability}

All analysis code, per-fold results, and figure sources accompany this article,
including the experiment scripts, the JSON reports they emit, and the
per-(model, fold) CSVs underlying every table and figure. The résumé--posting
corpus analysed is publicly available on the Hugging Face Hub; as
\S\ref{sec:ml} records, it carries no licence statement, and we analysed it as
distributed and redistribute none of it.

No human subjects were involved in the studies reported here. No résumé, job
posting, or candidate record was collected, contacted, or re-identified; the
face-verification study uses the public LFW benchmark~\cite{huang2007lfw}, and
the grounding evaluation uses a synthetic résumé corpus.

\section*{Declaration of Competing Interests}

The author declares no competing financial or non-financial interests. This work
received no external funding and was not commissioned, sponsored, or reviewed by
any vendor of hiring software.

\bibliographystyle{ACM-Reference-Format}
\bibliography{references}

\end{document}
